\chapter{Introduction}
\label{chap:introduction}
\section{Background}
Recent advancements in computer hardware and software have enabled nearly photorealistic and real-time
simulation of virtual worlds. Technologies such as augmented reality (AR) and virtual reality (VR) 
enhance the immersion of these simulated world by overlaying information over the user’s field of view.
Since the technology is still in its early stages, there is still a lot of unsolved problems and progress 
to be made. The immersion into virtual worlds is usually achieved through the use of a headset. Early
headsets were really bulky, had low resolution and short battery life. Naturally most of the research
was focused on tackling these problems. In the present day we have consumer grade devices like meta quest 3
and apple vision pro that offer excelent battery life, sharp visuals and immersive audio. In the
future we could expect this technology to shrink to the size of regular pair of glasses. Resarchers
at Meta are already stepping in that direction with ther recently unveiled Orion glasses. 

\bigskip \noindent
The hardware improvements also made the creation of Artificial Intelligence (AI) possible. It was theorized as early as 1950s by Alan Turing if machines are capable of thinking \cite{10.1093/mind/LIX.236.433}. In 1959 Arthur Samuel coined the term 'machine learning', but Ii was only in 1990s that hardware and algorithms were sufficiently advanced that machine learning began to gain traction.  Due to well-defined environments and rules, games were the first domain to notice
the impact of AI. In 1997, Deep Blue 
\cite{deepblue} managed to beat the world chess champion Garry
Kasparov. As computers following Moore’s law \cite{moore1998cramming} grew in power, more breakthroughs
began to emerge. Currently, in 2025, the AI is capable creating images, videos, music, text, research podcasts and even interactive worlds. 

\bigskip \noindent
The artificial intelligence also gave rise to the robotics industry. AI is transforming robots from rigid, pre-programmed machines into intelligent systems capable of performing increasingly complex tasks. Robots are often trained in simulated environments as this allows for faster and safer training without needing to worry about demaging robot or the environment.

\section{Problem statement}
Despite the fact that VR and AR technologies are rapidly improving, there is still a major problem,
these technologies usually immerse only two senses: sight and
hearing, but we, as humans, have five: sight, hearing, taste, smell, and touch. This means that there is still a lot of room for improvement as inclusion of additional senses would provide a more immersive VR/AR experience. 

\section{Objective}
The thesis aims to address
the limitation of VR technology by incorporating the sense of touch into the computer-simulated virtual worlds. The are two objectives, first, following engineering design process, will be focused on developing a portable exoskeleton capable of giving it's user an immersive sense of touch in VR. Second, will be dedicated to exploring if the developed exoskeleton could aid in the process of training robots in simulated world.

\section{Research question}
Can a portable exoskeleton that gives it's users the ability to touch virtual objects be created in a way that it is a viable solution to aid in the training of the robots in the simulated environments?